% =============================================================================
% Chapter 2 — Literature Review
% Tara Torbati — Master's Thesis
% =============================================================================

\chapter{Literature Review}
\label{chap:lit}

This chapter surveys the technical literature relevant to constrained
optimal irrigation control on topographically heterogeneous agricultural
fields. Section~\ref{sec:lit-context} situates the problem within the
broader landscape of agricultural water management. Sections
\ref{sec:lit-classical} through \ref{sec:lit-rl} review the principal
control architectures that have been proposed for irrigation, in
roughly increasing order of mathematical sophistication and
computational demand: classical and heuristic controllers, Model
Predictive Control with its modern variants, and Reinforcement
Learning. Section~\ref{sec:lit-spatial} addresses the specific question
of spatially distributed control and topographic flow modeling, which
is central to the present work but treated only obliquely in most of
the irrigation-control literature. Section~\ref{sec:lit-sensing}
briefly reviews the sensing and communication infrastructure that
supports field deployment. Section~\ref{sec:lit-deficit} comments on
deficit-irrigation strategies for completeness. Section~\ref{sec:lit-gap}
synthesizes the surveyed work and identifies the specific research
gap that this thesis addresses.

\section{Context: Agricultural Water Management Under Scarcity}
\label{sec:lit-context}

Agriculture is the largest consumer of freshwater globally, accounting
for approximately 70\% of human withdrawals~\cite{Sharma2025,Bayar2025}.
Under demographic pressure projected to push global food demand sharply
upward by mid-century, alongside the deepening effects of regional
water scarcity, the agricultural sector is undergoing a transition from
volume-driven, heuristically managed irrigation toward data-driven
precision systems that match water application to crop demand at
spatial and temporal scales finer than any human operator can
manually achieve~\cite{Sharma2025,LopezJimenez2022,Mgana2026}.

This transition is particularly acute in Iran, where surface water and
groundwater resources have been pushed beyond sustainable yield by a
combination of supply-side infrastructure expansion, subsidized energy
for groundwater pumping, and weak volumetric controls on agricultural
extraction~\cite{Mesgaran2018,Nouri2023,Behling2022}. In Gilan
Province, the principal rice-producing region, the combination of
heavily subsidized water tariffs and informal extraction practices has
removed economic constraints on use, with documented consequences for
land subsidence, stream desiccation, and aquifer
collapse~\cite{Mesgaran2018,PourgholamAmiji2024}. The resulting
governance asymmetry --- where price signals fail to discipline
demand and where conservation measures alone produce systemic rebound
effects --- has been argued to make purely policy-based
demand-management approaches insufficient, and to motivate technical
controls that operate on the field directly through water-budget
constraints~\cite{PerezBlanco2020,Hajirad2025,Hesabi2023}.

Within this technical-control framing, the irrigation problem becomes a
constrained finite-horizon optimal control problem: allocate a finite
volumetric water budget over a fixed agricultural season to maximize
crop yield (or some economic equivalent) while respecting per-day
actuator capacities, soil-physical constraints, and crop-physiological
tolerances. The remainder of this chapter reviews the control
architectures that have been proposed to solve this problem.

\section{Classical and Heuristic Controllers}
\label{sec:lit-classical}

The dominant irrigation control regime in deployed agricultural systems
remains heuristic: fixed schedules calibrated to long-run climatic
norms, threshold-based on/off control triggered by a single soil
moisture sensor, or open-loop time-based timers~\cite{Sharma2025,Bayar2025}.
These approaches are simple to implement, robust to sensor failure,
and operationally familiar; they are rightly the first option for
small-scale and resource-constrained deployments. They are also,
however, fundamentally reactive. They cannot anticipate impending
drought or rainfall, cannot exploit weather forecasts, and cannot
optimize against a global seasonal water budget --- the controller
sees only present state, not future demand.

A more sophisticated heuristic class is fuzzy-logic control, in which
linguistic rules over input ranges (e.g.\ ``IF soil moisture is LOW
AND temperature is HIGH THEN irrigation is HIGH'') are compiled into
crisp control actions through a defuzzification step. Fuzzy systems
have been integrated with PID feedback to handle the nonlinear,
delayed dynamics of soil-water response in fertigation systems and in
canal-gate control~\cite{Sharma2025,Mahmud2026}. Variants that
incorporate online self-correction have been demonstrated to improve
transient response in nonlinear nutrient regulation tasks. While
fuzzy-PID architectures partially address the nonlinearity limitation
of pure PID, they share with all rule-based methods the inability to
optimize systematically over time or to enforce strict resource
constraints --- a fuzzy controller has no native mechanism to
guarantee that cumulative seasonal water use stays below a fixed
budget.

These limitations become operationally consequential in water-scarce
contexts where the binding constraint is not the per-day soil
moisture target but the season-cumulative water allocation. In such
regimes, the central question shifts from ``how much water does the
crop need today'' to ``how should a finite seasonal water budget be
allocated across 90+ days of varying demand to maximize terminal
yield,'' a question that classical and heuristic controllers cannot
formally address. This motivates the transition to optimization-based
control, taken up in the next section.

\section{Optimal and Predictive Control}
\label{sec:lit-mpc}

\subsection{The Optimal Control Formulation of Irrigation}

The transition from heuristic to optimal control begins with the
formal recognition that irrigation is a constrained finite-horizon
optimal control problem. Following the formulation in
\cite{LopezJimenez2022,Lopezjimenez2024}, the discrete-time problem
is to determine a control sequence
$\mathbf{u}^*(\cdot) = \arg\min_{\mathbf{u}(\cdot)} J(\mathbf{x}(\cdot), \mathbf{u}(\cdot))$
subject to dynamics
$\mathbf{x}(k+1) = f(\mathbf{x}(k), \mathbf{u}(k), \mathbf{w}(k))$,
state constraints (soil moisture remains above the wilting point),
input constraints (per-day actuator capacity), and a season-cumulative
budget constraint
$\sum_k \sum_n u^{(n)}(k) \leq W_{\mathrm{total}}$.
The cost $J$ typically combines a terminal-state reward (terminal
biomass) with running penalties on water consumption, soil moisture
deficit, and surface ponding.

This formulation captures three features that heuristic controllers
cannot: (i) an explicit forecast horizon over which future weather and
crop development are anticipated; (ii) a global resource constraint
that couples decisions across the entire season; and (iii) a
multi-objective cost function with weights that admit economic
calibration. The optimization itself is generally a non-convex
nonlinear program due to crop-physiology nonlinearities and to the
saturation regime of soil drainage, requiring interior-point or
sequential-quadratic-programming methods for solution.

\subsection{Receding-Horizon Implementation: Model Predictive Control}

The standard approach to applying optimal control online is Model
Predictive Control (MPC). At each timestep, MPC solves the optimal
control problem over a finite prediction horizon $H_p$, applies only
the first-step control action, advances the state by one timestep,
and re-solves the problem with updated state and forecast information.
The receding-horizon mechanism transforms an open-loop optimization
into a feedback control law and provides a degree of robustness to
modeling errors and disturbances~\cite{LopezJimenez2022,Sharma2025}.

For irrigation, MPC has been demonstrated to improve water-use
efficiency by 8--29\% relative to heuristic and open-loop schedulers
in greenhouse and open-field studies~\cite{LopezJimenez2022,Mahmud2026}.
The key technical challenges are (i) constructing a control-oriented
model that is faithful enough to predict the principal hydrological
and biological responses while remaining computationally tractable
within the receding-horizon loop, (ii) handling the non-smoothness of
crop and soil dynamics (saturation thresholds, $\max$ and $\min$
operators) within smooth-gradient solvers, and (iii) calibrating the
multi-term cost function to reflect the agronomic and economic
priorities of the application context.

\subsection{Recent Variants and Architectural Extensions}

Several recent variants of irrigation MPC address specific limitations
of the standard formulation. \emph{Economic MPC} formulations replace
the conventional setpoint-tracking objective with a direct economic
cost that incorporates real-time price signals for water and energy,
enabling controllers to respond to time-varying tariff
structures~\cite{Sharma2025}. \emph{Rainfall-Intensity-and-Soil-Properties
MPC (RISPMPC)} extends standard MPC by integrating the Mein--Larson
infiltration model to anticipate runoff and percolation losses under
predicted rainfall, allowing the controller to deliberately undershoot
soil moisture targets in advance of an incoming storm so that the
soil matrix has capacity to absorb the rainfall productively. Reported
water savings against conventional MPC reach 30+\% on heavy-clay
soils, where infiltration capacity is the binding constraint~\cite{Sharma2025}.

Distributed MPC architectures address the scalability of the
optimization to large field domains and irrigation networks. Negenborn
et al.~\cite{negenborn2009distmpc} formalized distributed MPC for
multi-reach canal networks, decomposing a single large optimization
into per-subsystem optimizations coupled through inter-subsystem
communication. Subsequent work extended this paradigm to hierarchical
distributed MPC for basin-scale canal networks under water scarcity,
enabling tractable real-time optimization across irrigation districts
spanning thousands of hectares~\cite{LopezJimenez2022}. These
distributed approaches share with the present work the recognition
that single-monolithic-NLP MPC scales poorly with the number of
controlled subsystems.

Data-driven MPC variants substitute the physics-based predictive model
with learned surrogates, most prominently Long Short-Term Memory
(LSTM) networks trained on offline simulation data. These approaches
trade physical interpretability for adaptability across heterogeneous
field conditions, and they have been integrated with mixed-integer
formulations to handle discrete actuator levels in spatially variable
fields~\cite{Maniam2026}.

\section{Reinforcement Learning for Irrigation}
\label{sec:lit-rl}

\subsection{Why Reinforcement Learning}

The principal limitation of MPC is the per-step computational burden
of solving the nonlinear program online. Even with warm-starting and
smooth approximations, MPC solve times on field-scale problems range
from seconds at short horizons to minutes or hours at long
horizons~\cite{Lopezjimenez2024}. This places MPC out of reach for
deployment on edge computing devices, low-power embedded controllers,
or systems requiring sub-second control updates.

Reinforcement Learning (RL) addresses this latency limitation
structurally. After offline training in a simulated environment, an
RL policy is a fixed-parameter neural network whose forward-pass
inference is deterministic and bounded by the network's
fixed-architecture compute requirement --- typically sub-millisecond
on commodity hardware regardless of the underlying problem
dimensionality. The training cost is amortized over many millions of
environment interactions but pays off at every deployment step.

\subsection{Applications of RL to Irrigation Scheduling}

RL has been applied to irrigation across crop types, geographies, and
algorithm classes. Q-learning agents trained against the AquaCrop
crop-soil simulator have produced policies that outperform heuristic
schedules on lettuce production in Brazil~\cite{Sharma2025}. Deep
RL agents trained with high-dimensional state observations (soil
moisture across multiple depths, leaf-area index, accumulated rainfall
and irrigation, phenological stage) have been demonstrated to improve
profit in open-field wheat production by up to 17\% relative to
conventional heuristic rules, with policies that adapt to scenario
variation in ways that fixed rules cannot~\cite{Bayar2025,Hu2025}. In
controlled-environment agriculture, Advantage Actor-Critic (A2C)
agents have demonstrated water-use efficiency gains over both
threshold-based and time-based controllers across distinct growth
phases of greenhouse crops~\cite{Bayar2025}.

The principal practical challenge across these applications is
data-efficiency. Direct training on physical fields is infeasible
due to the cost of irreversibly damaging real crops during the
exploration phase. The dominant practice is therefore to train
against a high-fidelity simulator (such as AquaCrop, APSIM, or
custom agent-based models~\cite{Lopezjimenez2024}) and then transfer
the learned policy to the physical environment. The fidelity of the
simulator becomes the bottleneck: a poorly-validated simulator
produces policies that do not transfer.

\subsection{Continuous-Action Control: From DDPG to Soft Actor-Critic}

For irrigation, the action space is continuous: per-agent water
volume in some bounded interval. This rules out value-iteration
methods that operate on discrete action sets and requires
policy-gradient or actor-critic methods compatible with continuous
controls. The Deep Deterministic Policy Gradient (DDPG) algorithm of
Lillicrap et al.~\cite{lillicrap2015ddpg} was the first deep RL method
to demonstrate competitive performance on continuous control tasks,
combining off-policy Q-learning with a deterministic policy gradient
through an actor-critic architecture. DDPG, however, exhibited high
sensitivity to hyperparameters and instability across random seeds,
limiting its applicability outside well-tuned benchmark settings.
Twin Delayed Deep Deterministic Policy Gradient
(TD3)~\cite{fujimoto2018td3} addressed several of these failure modes
through twin Q-networks, delayed policy updates, and target-policy
smoothing, producing more stable training but retaining the
deterministic-policy framework.

The Soft Actor-Critic (SAC) algorithm of Haarnoja et al.\
\cite{haarnoja2018sac,haarnoja2019sac} is the current default for
continuous control in deep RL. SAC is built on the maximum-entropy
reinforcement learning framework, in which the standard objective
$\mathbb{E}_\pi[\sum_t r_t]$ is augmented with an entropy bonus
$\mathcal{H}[\pi(\cdot|s_t)]$ that rewards stochastic exploration:
\begin{equation}
J(\pi) = \mathbb{E}_\pi\Bigl[\sum_{t=0}^{\infty} \gamma^t
(r_t + \alpha \mathcal{H}[\pi(\cdot \mid s_t)])\Bigr]
\label{eq:sac-objective}
\end{equation}
The temperature parameter $\alpha$ trades off reward maximization
against entropy maximization and is automatically tuned in the modern
formulation~\cite{haarnoja2019sac} via a constrained Lagrangian
update against a target entropy. The maximum-entropy framework has
two principal practical benefits relative to deterministic-policy
methods: it produces stochastic policies that explore more
effectively during training, and it produces policies that capture
multi-modal optima --- a property of value when several distinct
control sequences achieve similar reward, as is common in irrigation
where the timing of water application admits considerable latitude.
SAC is off-policy (using a replay buffer) and uses twin Q-networks
in the manner of TD3 to mitigate value overestimation.

For the irrigation problem of this thesis, SAC is adopted as the RL
architecture. The justification is threefold: continuous action
support without quantization artifacts, sample efficiency on the
order required for high-fidelity simulator training, and demonstrated
robustness across hyperparameter and seed variation in the maximum-
entropy formulation~\cite{haarnoja2019sac}.

\subsection{Reward Engineering and Constraint Handling}

A central challenge in applying RL to constrained problems is that
standard RL algorithms have no native mechanism for enforcing hard
constraints~\cite{Bayar2025}. The MPC solver, by contrast, treats the
seasonal water budget as a hard linear inequality enforced at every
solver iteration. RL alternatives rely on reward shaping (penalizing
constraint violations) and on action clipping (truncating actions to
the feasible set), but these mechanisms shift constraint satisfaction
from a guarantee to an emergent property of training. The reward
function therefore requires careful construction. Documented practice
in irrigation RL combines: a positive baseline reward for matching
crop water demand, an exponential penalty for over-irrigation that
discourages dumping the budget early, a soft penalty on field-capacity
overshoot that prevents waterlogging stress, and a sparse terminal
bonus on harvested biomass that incentivizes long-horizon
yield~\cite{Bayar2025}. The cost function of the present thesis
follows this pattern, with an explicit policy-equivalence design
choice that the SAC reward is constructed as the exact negation of
the MPC path cost (see Chapter~\ref{chap:controllers}).

\subsection{Multi-Agent and Distributed Reinforcement Learning}

When the controlled system consists of many spatially distributed
sub-units, monolithic single-agent RL scales poorly because the
joint action space grows exponentially with the number of agents.
Multi-agent RL (MARL) addresses this by decomposing the problem into
per-agent policies that share information during training but execute
locally. The Centralized Training, Decentralized Execution (CTDE)
paradigm has emerged as the dominant
architecture~\cite{amato2024ctde}: during training, agents have
access to global state and all other agents' actions through a
centralized critic, while at deployment each agent acts on local
observations only. CTDE methods preserve the coordination benefits
of centralized training while producing policies that are scalable,
communication-free at deployment, and consistent with the partially
observable distributed structure of physical multi-agent systems.

For irrigation on a topographically heterogeneous field, the CTDE
paradigm is naturally applicable: each agent (corresponding to a
spatial cell of the field) observes its local state, while the
training-time critic has access to the joint state and can resolve
the credit-assignment problem of attributing global yield to per-cell
irrigation decisions. The thesis adopts this architecture; the
implementation details are presented in Chapter~\ref{chap:controllers}.

\section{Spatial Modeling and Topographic Heterogeneity}
\label{sec:lit-spatial}

The vast majority of irrigation control literature, including the
references reviewed in Sections~\ref{sec:lit-classical} through
\ref{sec:lit-rl}, treats the field as a single lumped soil-moisture
reservoir. This simplification is justified for small, flat,
homogeneous plots but breaks down when the controlled field is
spatially heterogeneous in elevation, soil texture, or crop
distribution. For agricultural fields with non-trivial topography ---
including the great majority of fields in mountainous regions such as
Gilan --- lateral water redistribution through surface runoff couples
the moisture states of adjacent cells in ways that lumped models
cannot represent.

\subsection{Variable-Rate Irrigation as a Spatial Control Precedent}

The closest commercial precedent for spatially distributed irrigation
control is Variable-Rate Irrigation (VRI), in which the application
rate of a center-pivot or linear-move sprinkler is modulated as the
machine traverses the field, applying different volumes to different
cells as a function of soil texture, slope, or
crop-vigor maps~\cite{Sharma2025,Bayar2025}. VRI partially addresses
the spatial-heterogeneity problem at the application end but does not
typically incorporate runoff coupling between cells: each cell's
prescription is computed independently from a static spatial map.
This is a reasonable approximation when slopes are mild and rainfall
events are gentle, but it loses fidelity in topographically active
fields where overland flow between cells is non-negligible.

\subsection{DEM-Based Flow Routing}

The hydrology literature has developed sophisticated machinery for
representing surface water flow on a discretized terrain. The
foundational technique is the D8 algorithm of O'Callaghan and
Mark~\cite{ocallaghan1984d8}, which assigns each cell of a digital
elevation model (DEM) a single downslope flow direction toward the
neighbor with the steepest descent. The D8 representation produces a
directed graph over the field cells whose edges encode lateral water
exchange. Multiple-flow-direction extensions, including
MD8~\cite{quinn1991mfd}, distribute flow proportionally among all
downslope neighbors weighted by slope; these reduce the artificial
flow-line concentration of D8 at the cost of greater computational
complexity. For control applications, where the routing graph must be
embedded in a symbolic optimization framework with second-order
gradient information, the simpler D8 representation has been the
practical choice in recent agent-based irrigation
work~\cite{Lopezjimenez2024,LopezJimenez2022}.

The integration of DEM-based routing into a control framework
introduces two design considerations not present in standard hydrology
applications: (i) the routing must be computable at every receding-
horizon iteration and must be differentiable with respect to control
inputs and state variables, ruling out iterative routing solvers and
discontinuous flow-direction switches; and (ii) the routing must
respect mass conservation across the field domain, including at the
field boundary, where naive routing schemes produce degenerate
sink behavior in low-elevation cells with no internal downhill
neighbors. The first consideration is addressed by static topological
sorting of the routing graph at NLP construction time, allowing the
forward sweep to be unrolled symbolically; the second is addressed
through fractional routing through padded boundary cells, as
implemented in the present thesis.

\subsection{Agent-Based Models for Spatial Irrigation Control}

The most recent line of work that directly addresses spatially
distributed irrigation control on heterogeneous fields is the
agent-based model of Lopez-Jimenez et al.~\cite{Lopezjimenez2024,LopezJimenez2022}.
This work formulates a discrete-cell model in which each agent
maintains its own soil moisture, ponding, and biomass states, coupled
through D8-derived surface runoff routing, and embeds this model
within a nonlinear MPC formulation for water-budget-constrained
irrigation. The agent-based formulation is compact relative to fine-
resolution finite-element soil hydrology and admits the symbolic
differentiation required by interior-point optimization, while
preserving the spatial fidelity needed to capture topographically
driven phenomena such as ridge runoff and valley waterlogging.

The present thesis builds directly on the Lopez-Jimenez framework with
three modifications motivated by the constrained control setting: an
explicit two-layer hydrological decoupling between surface ponding
and root-zone moisture, a padded-DEM boundary condition that resolves
a degenerate sink behavior in the original formulation, and a coupling
of drought stress directly to biomass accumulation rather than only to
maturation. These modifications and their motivations are documented
in Chapter~\ref{chap:system}.

\section{Sensing, Data Assimilation, and Communication Infrastructure}
\label{sec:lit-sensing}

Field deployment of any closed-loop irrigation controller depends on
the upstream sensing infrastructure that provides the soil moisture,
temperature, and weather observations consumed by the controller.
Soil-based sensing using Time Domain Reflectometry (TDR), Frequency
Domain Reflectometry (FDR), and capacitive probes is the dominant
modality, deployed at cell-level density across the
field~\cite{Sharma2025,Bayar2025}. Plant-based sensing of leaf
turgor, sap flow, and stomatal conductance is increasingly used in
high-value horticulture for direct measurement of physiological water
status. Atmospheric sensing using satellite remote sensing products
(SMAP, SMOS, MERRA-2 GWETROOT~\cite{Afrasiabikia2026}) provides
field-mean estimates that complement the localized in-situ
measurements and are particularly useful for validation of integrated
field-scale models.

The principal practical challenges are sensor drift in harsh field
environments, which motivates online data assimilation through
particle-filter or Kalman-filter calibration~\cite{Maniam2026}, and
the energy and bandwidth constraints of farm-wide deployment, which
motivates Low-Power Wide-Area Network (LPWAN) protocols such as LoRa
and NB-IoT for telemetry. These infrastructure questions, while
critical for deployment, are largely orthogonal to the control-policy
question of this thesis and are not developed further.

\section{Deficit Irrigation and Crop-Physiological Constraints}
\label{sec:lit-deficit}

A substantial literature addresses deficit-irrigation strategies in
which water application is deliberately reduced below crop demand
during non-critical phenological stages, exploiting the plant's
physiological response to mild water stress to improve fruit quality
and water-use efficiency. Regulated Deficit Irrigation (RDI) and
Partial Root-zone Drying (PRD) have been demonstrated to achieve
substantial water savings in stone fruits, processing tomatoes, and
citrus orchards while preserving or improving harvested
quality~\cite{Sharma2025}. These strategies are compelling for crops
in which water stress modulates harvest quality through secondary-
metabolite accumulation. They are, however, largely inapplicable to
flooded-paddy rice cultivation, where rice's intolerance of
moderate-to-severe drought stress~\cite{jalali2021,raoufi2025}
constrains the controller to keeping soil moisture above the
agronomic stress threshold throughout the season. The present thesis
operates in this constrained regime and does not pursue
deficit-irrigation strategies; the deficit budget scenarios at 85\%
and 70\% allocation evaluated in Chapter~\ref{chap:results} test
controller robustness under hard volumetric scarcity, not under
deliberate stress imposition.

\section{Synthesis and Research Gap}
\label{sec:lit-gap}

The literature surveyed in this chapter establishes the principal
families of control architectures that have been brought to bear on
the irrigation problem and their respective strengths and limitations.
Heuristic and classical controllers are simple and operationally
robust but cannot enforce season-cumulative constraints or anticipate
future weather. MPC provides constraint-aware, forecast-driven
optimization but at a per-step computational cost that limits
edge-device deployment and that scales unfavorably with prediction
horizon. RL provides sub-millisecond inference after training but
lacks native constraint guarantees and depends critically on the
fidelity of the training simulator. Each family has been demonstrated
in irrigation applications, and recent work has produced increasingly
sophisticated extensions (RISPMPC, distributed MPC, CTDE-based
multi-agent RL).

Within this landscape, several gaps motivate the present thesis.
First, the integration of agent-based topographic cascade routing
with constrained MPC has been demonstrated only in a handful of
recent works~\cite{Lopezjimenez2024,LopezJimenez2022}, and several
pathologies of the resulting NLP --- including degenerate sink
behavior, field-capacity overshoot at long horizons, and chronic
waterlogging in wet-year scenarios --- have not been systematically
characterized or resolved. Second, the calibration of the multi-term
MPC cost function to real economic price tiers has not been
demonstrated in the published literature; existing work treats the
weights as abstract regularization constants rather than as
quantities anchored to volumetric water tariffs. Third, the direct
head-to-head comparison of constrained MPC against a SAC controller
trained on the exact same plant model and against an exact-equivalent
reward function --- enabling rigorous policy-comparison rather than
algorithm-vs-algorithm benchmarking --- has not been reported for
spatially distributed agricultural control.

The thesis addresses these three gaps. Chapter~\ref{chap:system}
presents the validated agent-based virtual plant on which all
controllers operate. Chapter~\ref{chap:controllers} introduces a
constrained MPC formulation with a six-term cost function (later
reduced to five terms via factorial sensitivity analysis) anchored to
the four-tier Iranian water pricing schedule, and a SAC controller
trained against the exact-negation of the MPC path cost.
Chapter~\ref{chap:results} reports the comparative performance of
both controllers across nine combinations of climate scenario and
budget allocation, identifying the operational regimes in which each
architecture excels and quantifying the cost-function weight
calibration through a thirty-one-configuration sensitivity sweep.
The Iranian context provides both the policy motivation, through the
water-scarcity dynamics surveyed in Section~\ref{sec:lit-context},
and the empirical anchor, through the Gilan rice-production setting
on which the virtual plant is parameterized.
